\chapter{The 28-Dimensional Aspects: Geometric Foundation of Force Hierarchy}

% [The 28-Dimensional Aspects: Geometric Foundation of Force Hierarchy]

%----------------------------------------------------------------------
\section{Introduction}
%----------------------------------------------------------------------

\subsection{The Blind Spot Below~$c$}

Standard physics has no predictive structural model of the nucleus.  The strong force is modelled by virtual gluon exchange; nuclear structure is a ``drop-point'' composite of entities that create themselves, perform a function, and vanish.  This gives behaviours but no geometry.

SDT proposes that matter \emph{is} a superluminal toroidal vortex (a $6\pi$ trefoil knot rotating at ${\sim}1.836\,c$) confined within the $c$-boundary---a region invisible to direct observation.  The 28-dimensional aspect space describes the complete external interaction of such a vortex with the spation medium, without requiring access to the sub-$c$ interior.

\subsection{Scope}

This chapter:
\begin{itemize}
  \item defines the seven-level aspect hierarchy and each of the 28~aspects,
  \item derives the occlusion function $E(\mathbf{x},\hat{\mathbf{n}})$ and the $1/r^2$ force law,
  \item shows how the force hierarchy emerges from aspect-dependent coupling,
  \item formalises the $\Phi$-state dynamics and the movement principle,
  \item proposes three experimental tests.
\end{itemize}


%----------------------------------------------------------------------
\section{The Core Mechanism}
%----------------------------------------------------------------------

A single sentence encapsulates the physics:

\begin{quote}
\emph{Every particle occludes CMB pressure from reaching nearby spations, creating a pressure deficit; spations press against the particle to fill the deficit, and that pressure gradient \textbf{is} force.}
\end{quote}

\noindent Three facts constrain the mechanism:
\begin{enumerate}
  \item \textbf{Matter occludes CMB from spations.}  Spations do \emph{not} occlude each other; they flow inviscidly.
  \item \textbf{Pressure propagates at $c$} along zero-point lines from each particle.
  \item \textbf{The occlusion solid angle} $\Omega_{\text{blocked}}$ determines the coupling strength at every scale.
\end{enumerate}


%----------------------------------------------------------------------
\section{The 28-Dimensional State Manifold}
%----------------------------------------------------------------------

\subsection{Triangular-Number Architecture}

The complete state vector $\Xi \in \mathbb{R}^{28}$ is organised into seven levels, where Level~$N$ contains exactly $N$~aspects.  The total is the seventh triangular number:
\begin{equation}
  T_7 = \sum_{N=1}^{7} N = \frac{7 \times 8}{2} = 28.
\end{equation}
The number 28 is also the second perfect number ($1{+}2{+}4{+}7{+}14=28$), the number of pairwise interactions among 8~objects ($\binom{8}{2}=28$), and the dimension of the Lie algebra $\mathfrak{so}(8)$.


\subsection{Level~1: Zero-Point (1~aspect)}

\begin{center}
\begin{tabular}{clcl}
\hline
\# & Symbol & Units & Meaning \\
\hline
1 & $\xi_0$ & --- & Existence flag ($0$ or $1$) \\
\hline
\end{tabular}
\end{center}

If $\xi_0 = 1$ the particle exists and blocks CMB pressure along its zero-point line.  All subsequent aspects are gated by $\xi_0$; if $\xi_0 = 0$ the state vector is null.


\subsection{Level~2: Line (2~aspects)}

\begin{center}
\begin{tabular}{clcl}
\hline
\# & Symbol & Units & Meaning \\
\hline
2 & $\xi_{10}$ & m & Location (position vector) \\
3 & $\xi_{11}$ & m/s & Relocation (velocity vector) \\
\hline
\end{tabular}
\end{center}

Position and linear motion.  The zero-point line extends from $\xi_{10}$ and the pressure-wave arrival time at distance $r$ is
\begin{equation}
  t_{\text{arrival}} = t_0 + \frac{|\mathbf{r} - \xi_{10}|}{c}\,.
\end{equation}
For hydrogen at $r = a_0$:
\begin{equation}
  \frac{t_{\text{arrival}}}{T_{\text{orbital}}}
    = \frac{a_0/c}{2\pi a_0/v_1}
    = \frac{v_1}{2\pi c}
    = \frac{\alpha}{2\pi}
    = 1.162 \times 10^{-3}.
\end{equation}
Pressure is effectively instantaneous on the orbital timescale.


\subsection{Level~3: Plane (3~aspects)}

\begin{center}
\begin{tabular}{clcl}
\hline
\# & Symbol & Units & Meaning \\
\hline
4 & $\xi_{p0}$ & --- & Internal existence (planar boundary) \\
5 & $\xi_{p1}$ & m$^{2}$ & Planar relocation \\
6 & $\xi_{p2}$ & rad & Planar rotation \\
\hline
\end{tabular}
\end{center}

The 2D cross-section that determines directional occlusion.  For a linear molecule (e.g.\ N$_2$, bond length $1.10\times10^{-10}$\,m) the face-on to edge-on cross-section ratio is $\sigma_{\max}/\sigma_{\min} \approx 1.07$: even strongly non-spherical molecules are nearly isotropic occluders.


\subsection{Level~4: Sphere (4~aspects)}

\begin{center}
\begin{tabular}{clcl}
\hline
\# & Symbol & Units & Meaning \\
\hline
7  & $\xi_{s0}$ & m$^{3}$ & Shell existence (displaced volume) \\
8  & $\xi_{s1}$ & m$^{3}$/s & Shell relocation \\
9  & $\xi_{s2}$ & rad/s & Shell rotation \\
10 & $\xi_{s3}$ & --- & Orientation (rotation-axis unit vector) \\
\hline
\end{tabular}
\end{center}

The effective radius is $R_{\text{eff}} = (3\,\xi_{s0}/4\pi)^{1/3}$ and the solid-angle occlusion at distance $r \gg R$ is
\begin{equation}\label{eq:E-sphere}
  E = \frac{\Omega_{\text{blocked}}}{4\pi} \approx \frac{R_{\text{eff}}^2}{4\,r^2}\,.
\end{equation}
This is the origin of the $1/r^2$ force law: the occlusion solid angle subtended by a sphere falls as $1/r^2$, so the resulting pressure gradient---and therefore force---scales identically.


\subsection{Level~5: Torus (5~aspects)---Matter Structure}

\begin{center}
\begin{tabular}{clcl}
\hline
\# & Symbol & Units & Meaning \\
\hline
11 & $T_1$ & m & Central ring (zero-point line, max compression) \\
12 & $T_2$ & m & Tube diameter (vortex thickness) \\
13 & $T_3$ & m$^{2}$ & Topological surface (torus surface area) \\
14 & $T_4$ & m$^{3}$\,Pa & Polarised volume (aperture $\times$ pressure gradient) \\
15 & $T_5$ & Pa/m & Aspect gradation (internal pressure gradient) \\
\hline
\end{tabular}
\end{center}

Matter \emph{is} a toroidal vortex.  The five torus aspects encode its geometry:
\begin{itemize}
  \item \textbf{$T_1$ is the zero-point line}---the constriction line along which pressure propagates at~$c$.  Every particle has $T_1 > 0$.
  \item \textbf{$T_3$ determines occlusion:}
    \begin{equation}\label{eq:E-torus}
      E = \frac{T_3}{4\pi\,r^2}\,.
    \end{equation}
  \item $T_4$ generates compression (gravitational waves).
  \item $T_5$ determines self-screening.
\end{itemize}

\noindent The torus surface area is $T_3 = \pi^2\,T_1\,T_2$ and the torus volume is $V = \tfrac{1}{2}\pi^2\,T_1\,T_2^2$.

This level is why particles have spin (real rotation), magnetic moments (circulating current), and binding energy (energy required to break the torus topology).


\subsection{Level~6: Dynamism (6~aspects)---Time Evolution}

\begin{center}
\begin{tabular}{clcl}
\hline
\# & Symbol & Units & Meaning \\
\hline
16 & $\Phi_0$ & sr & Omnidirectionality (typically $4\pi$) \\
17 & $\Phi_1$ & m/s$^{2}$ & Dynamic translocation \\
18 & $\Phi_2$ & Hz & Oscillation (precession, reversal) \\
19 & $\Phi_3$ & $\pm1$ & Inversion / chirality \\
20 & $\Phi_4$ & --- & State trajectory variance \\
21 & $\Phi_5$ & J & Phase-transition threshold \\
\hline
\end{tabular}
\end{center}

\subsubsection{$\Phi$-State Dynamics}

A system configuration $Z$ occupies one of four flow states:
\begin{itemize}
  \item $\Phi_1$: Free-flow (no confinement, no circulation).
  \item $\Phi_2$: Boundary-locked (stable confinement, minimal circulation).
  \item $\Phi_3$: Toroidal circulation (steady helical wake).
  \item $\Phi_4$: Shunt-dominated (exchange with external reservoirs).
\end{itemize}

Open systems transition $\Phi_2 \to \Phi_3 \to \Phi_4$ under external pressure gradients.  Closed systems stabilise at $\Phi_2$ or $\Phi_3$.

For hydrogen, the $\Phi_2 \to \Phi_3$ transition corresponds to the first excitation:
\begin{equation}
  E(\Phi_2 \to \Phi_3) = E_2 - E_1 = 10.205\;\text{eV} = \text{Lyman-}\alpha\text{ energy.}
\end{equation}


\subsection{Level~7: Energy (7~aspects)---Force Manifestation}

\begin{center}
\begin{tabular}{clcl}
\hline
\# & Symbol & Units & Meaning \\
\hline
22 & $\varepsilon_0$ & J & Potential energy \\
23 & $\varepsilon_1$ & J & Kinetic energy (bulk motion) \\
24 & $\varepsilon_2$ & J & Rotational energy \\
25 & $\varepsilon_3$ & J & Field energy (pressure-occlusion field) \\
26 & $\varepsilon_b$ & J & Binding energy (decoherence threshold) \\
27 & $\varepsilon_4$ & W & Flux (energy-transfer rate, propagates at $c$) \\
28 & $\varepsilon_5$ & J & Transmission (sound, percussion, gravitational waves) \\
\hline
\end{tabular}
\end{center}

\begin{itemize}
  \item $\varepsilon_3$ is the energy stored in the pressure deficit when matter blocks CMB from spations.
  \item $\varepsilon_4$ is the flux: pressure waves propagating at~$c$.
  \item $\varepsilon_5$ includes \textbf{gravitational waves}, which SDT identifies as compression waves in the spation lattice.
\end{itemize}


%----------------------------------------------------------------------
\section{Minimality}
%----------------------------------------------------------------------

Dropping any single level removes a necessary degree of freedom:

\begin{center}
\begin{tabular}{lll}
\hline
\textbf{If removed} & \textbf{Cannot describe} & \textbf{Consequence} \\
\hline
Level~1 (Point) & creation/annihilation & no particle counting \\
Level~2 (Line) & position, motion & no trajectories \\
Level~3 (Plane) & cross-section shape & no directional dependence \\
Level~4 (Sphere) & volume, orientation & no $1/r^2$ law \\
Level~5 (Torus) & matter topology & no spin, no binding \\
Level~6 (Dynamism) & time evolution & no dynamics \\
Level~7 (Energy) & force magnitudes & no quantitative predictions \\
\hline
\end{tabular}
\end{center}

\noindent The 28-aspect space is therefore a \textbf{minimal complete basis} for SDT dynamics.


%----------------------------------------------------------------------
\section{Occlusion Function and the Force Law}
%----------------------------------------------------------------------

\subsection{Definition}

The occlusion fraction at location $\mathbf{x}$ along direction $\hat{\mathbf{n}}$ is
\begin{equation}
  E(\mathbf{x},\hat{\mathbf{n}}) = \frac{\Omega_{\text{blocked}}(\mathbf{x},\hat{\mathbf{n}})}{4\pi}\,,
  \quad E \in [0,1].
\end{equation}
$E = 0$: no occlusion.  $E \to 1$: near-total occlusion.

In terms of the torus aspects:
\begin{equation}
  E = \frac{T_3}{4\pi\,r^2}\,,
\end{equation}
where $r$ is the distance from the particle.

\subsection{Coupling Factor}

The effective coupling factor is
\begin{equation}
  \mathcal{C}(\mathbf{x},\hat{\mathbf{n}}) = 1 - E(\mathbf{x},\hat{\mathbf{n}}).
\end{equation}

\subsection{General Force}

The interaction between two displacements with compactnesses $\kappa_1, \kappa_2$ at separation~$r$:
\begin{equation}\label{eq:force}
  F(\kappa_1, \kappa_2, r)
    = \frac{\kappa_1\,\kappa_2}{4\pi\,K_{\text{bulk}}\,r^2}\;\mathcal{C}(\mathbf{x},\hat{\mathbf{n}})\,,
\end{equation}
where $K_{\text{bulk}}$ is the spation bulk modulus.  Different force regimes correspond to different \textbf{aspect-controlled limits} of $\mathcal{C}$ and $\kappa$.

\subsection{Pressure Wave Propagation}

From each particle, the pressure field is
\begin{equation}
  P(r, t) = P_{\text{CMB}} - \Delta P(r) \times H\!\left(t - t_0 - \frac{r}{c}\right),
\end{equation}
where $H(t)$ is the Heaviside step function (the wave has not arrived if $t < t_0 + r/c$) and the pressure deficit is
\begin{equation}
  \Delta P(r) = P_{\text{CMB}} \times \frac{T_3}{4\pi\,r^2}\,.
\end{equation}
The spation-side force density is then
\begin{equation}
  \mathbf{F} = -V_{\text{disp}}\,\nabla P
    = V_{\text{disp}}\,P_{\text{CMB}}\,\nabla E\,.
\end{equation}


%----------------------------------------------------------------------
\section{Force Hierarchy from Aspect Coupling}
%----------------------------------------------------------------------

SDT predicts a hierarchy of effective screening factors $S = E$:
\begin{equation}
  S_{\text{nuclear}} \sim 1,
  \quad
  S_{\text{atomic}} \sim 10^{-9},
  \quad
  S_{\text{planetary}} \sim 10^{-30},
  \quad
  S_{\text{cosmological}} \sim 10^{-123}.
\end{equation}

This single mechanism spans the electromagnetic-to-gravitational ratio and resolves the hierarchy problem \textbf{without introducing new fields or particles}.

\begin{center}
\begin{tabular}{llll}
\hline
\textbf{Scale} & \textbf{$S$} & \textbf{Regime} & \textbf{Dominant aspects} \\
\hline
Nuclear (${\sim}$fm) & $\sim 1$ & Strong & $T_1, T_2, T_3$ (torus overlap) \\
Atomic (${\sim}$pm) & $\sim 10^{-9}$ & Electromagnetic & $\xi_{s0}, T_3$ (shell occlusion) \\
Planetary (${\sim}$AU) & $\sim 10^{-30}$ & Gravity & $\xi_{10}, \varepsilon_3$ (cumulative field) \\
Cosmological & $\sim 10^{-123}$ & Dark energy / $\Lambda$ & $\Phi_0$ (omnidirectional CMB) \\
\hline
\end{tabular}
\end{center}


%----------------------------------------------------------------------
\section{Movement Principle and Choice Gradient}
%----------------------------------------------------------------------

\subsection{Principle}

Systems evolve toward configurations that maximise interaction options:
\begin{equation}
  \frac{dN_{\text{choices}}}{dt} \ge 0\,,
\end{equation}
where $N_{\text{choices}}$ counts distinct coupling configurations allowed by the 28~aspects.

\subsection{Choice Potential}

Define
\begin{equation}
  \Psi = \ln N_{\text{choices}}\,.
\end{equation}
The movement principle becomes
\begin{equation}
  \frac{d\Psi}{dt}
    = \nabla_{\text{aspect}}\Psi \cdot \dot{\mathbf{A}} \ge 0\,,
\end{equation}
where $\mathbf{A} \in \mathbb{R}^{28}$ is the full aspect vector.  This yields deterministic drift toward higher-connectivity states.

The movement principle is SDT's geometric analogue of the Second Law of Thermodynamics: systems do not maximise entropy; they maximise \emph{interaction choices}.


%----------------------------------------------------------------------
\section{Gravitational Waves as Compression Waves}
%----------------------------------------------------------------------

In this framework gravitational waves are compression waves in the spation lattice, encoded by three aspects:

\begin{itemize}
  \item $T_4$ (polarised volume): creates the compression.
  \item $\Phi_2$ (oscillation): creates the wave frequency.
  \item $\varepsilon_5$ (transmission): carries the wave energy.
\end{itemize}

Formally:
\begin{equation}
  h_{\mu\nu}(t, r)
    = \sum_{\text{particles}} \bigl[T_4 \times \Phi_2 \times \delta\!\left(t - r/c\right)\bigr].
\end{equation}


%----------------------------------------------------------------------
\section{Experimental Tests}
%----------------------------------------------------------------------

\subsection{Test~1: Anisotropy Detection}

\textbf{Setup:} Resonant cavity oriented along different spatial directions.

\textbf{Prediction:} Coupling variations correlate with the occlusion direction $\hat{\mathbf{n}}$.  At Earth's surface, $E_{\downarrow} \approx \pi/(4\pi) = 0.25$ (hemisphere blocked) while $E_{\uparrow} \approx 0$ (open sky).  Mode frequencies in a vertical cavity should shift by $\sim 10^{-15}$ relative to a horizontal cavity.

\textbf{Tests aspects:} $\xi_{s3}$ (orientation), $\xi_{p2}$ (planar rotation).

\subsection{Test~2: $\Phi$-State Identification}

\textbf{Setup:} Controlled vortex system (superfluid helium, BEC).

\textbf{Prediction:} At the $\Phi_2 \to \Phi_3$ transition, helical wake signatures appear.  In superfluid He-4 the circulation quantum is $\kappa = h/m_{\text{He}} = 9.97 \times 10^{-8}\;\text{m}^2\text{/s}$, which SDT interprets as the minimum $T_1 \times v_1$ product for a stable toroidal vortex.

\textbf{Tests aspects:} $\Phi_2$ (oscillation), $\Phi_3$ (chirality), $T_1$ (central ring).

\subsection{Test~3: Choice-Gradient Verification}

\textbf{Setup:} Multi-vortex assembly (vortex lattice in rotating superfluid).

\textbf{Prediction:} Vortices migrate toward configurations with more interaction paths---i.e.\ hexagonal lattice (6~neighbours) rather than square (4~neighbours).  This is directly observed experimentally.

\textbf{Tests aspects:} $\Psi$ (choice potential), $d\Psi/dt \ge 0$ (movement principle).


%----------------------------------------------------------------------
\section{Summary of the Complete Aspect Table}
%----------------------------------------------------------------------

\begin{table}[h]
\centering
\small
\begin{tabular}{|c|c|c|c|p{4.0cm}|p{4.0cm}|}
\hline
\textbf{Lvl} & \textbf{\#} & \textbf{Symbol} & \textbf{Units} & \textbf{Physical meaning} & \textbf{Role in pressure blocking} \\
\hline
1 & 1  & $\xi_0$       & ---        & Existence                  & Particle exists $\to$ blocks \\
\hline
2 & 2  & $\xi_{10}$    & m          & Position                   & Zero-point line origin \\
2 & 3  & $\xi_{11}$    & m/s        & Velocity                   & Doppler correction \\
\hline
3 & 4  & $\xi_{p0}$    & ---        & Planar boundary            & 2D cross-section \\
3 & 5  & $\xi_{p1}$    & m$^2$      & Planar position            & Orientation within plane \\
3 & 6  & $\xi_{p2}$    & rad        & Planar rotation            & Blocking directionality \\
\hline
4 & 7  & $\xi_{s0}$    & m$^3$      & Volume                     & Effective radius $R$ \\
4 & 8  & $\xi_{s1}$    & m$^3$/s    & Volume change              & Expansion/contraction \\
4 & 9  & $\xi_{s2}$    & rad/s      & Shell rotation             & Angular momentum \\
4 & 10 & $\xi_{s3}$    & ---        & Orientation                & Zero-point line direction \\
\hline
5 & 11 & $T_1$         & m          & \textbf{Central ring}      & \textbf{Zero-point line ($v = c$)} \\
5 & 12 & $T_2$         & m          & Tube diameter              & Displacement volume \\
5 & 13 & $T_3$         & m$^2$      & \textbf{Surface area}      & \textbf{$E = T_3/(4\pi r^2)$} \\
5 & 14 & $T_4$         & m$^3$\,Pa  & Polarised volume           & Compression waves \\
5 & 15 & $T_5$         & Pa/m       & Aspect gradation           & Self-screening \\
\hline
6 & 16 & $\Phi_0$      & sr         & Omnidirectionality         & Full-sphere blocking \\
6 & 17 & $\Phi_1$      & m/s$^2$    & Acceleration               & Higher-order motion \\
6 & 18 & $\Phi_2$      & Hz         & Oscillation                & Wave frequency \\
6 & 19 & $\Phi_3$      & $\pm 1$    & Chirality                  & Handedness \\
6 & 20 & $\Phi_4$      & ---        & State variance             & External response \\
6 & 21 & $\Phi_5$      & J          & Phase transition           & Structural change \\
\hline
7 & 22 & $\varepsilon_0$ & J        & Potential                  & Position energy \\
7 & 23 & $\varepsilon_1$ & J        & Kinetic                    & Motion energy \\
7 & 24 & $\varepsilon_2$ & J        & Rotational                 & Spin energy \\
7 & 25 & $\varepsilon_3$ & J        & \textbf{Field}             & \textbf{Pressure-occlusion field} \\
7 & 26 & $\varepsilon_b$ & J        & Binding                    & Structure stability \\
7 & 27 & $\varepsilon_4$ & W        & \textbf{Flux}              & \textbf{Propagates at $c$} \\
7 & 28 & $\varepsilon_5$ & J        & \textbf{Transmission}      & \textbf{Grav.\ waves = compression} \\
\hline
\end{tabular}
\caption{Complete 28-aspect state manifold ($\Xi \in \mathbb{R}^{28}$).  Bold entries are the three aspects most critical to the pressure-blocking mechanism.}
\end{table}


%----------------------------------------------------------------------
\section{Discussion}
%----------------------------------------------------------------------

The seven-level hierarchy unifies SDT's geometric primitives with force scaling, screening, and system evolution.  It avoids new particle species and attributes the force hierarchy entirely to occlusion-controlled coupling and flow-state constraints.

The triangular structure $1{+}2{+}3{+}4{+}5{+}6{+}7 = 28$ is not imposed; it emerges from the progression Point $\to$ Line $\to$ Plane $\to$ Sphere $\to$ Torus $\to$ Dynamism $\to$ Energy, each level inheriting and extending the previous.  The fact that $28 = \binom{8}{2}$ (the pairwise interaction count of 8~objects) and $28$ is a perfect number may be geometrically significant; this connection deserves further investigation.

All standard-model observables---spectral lines, binding energies, magnetic moments, scattering cross-sections---are in principle computable from the 28-aspect state vector and the occlusion integral $E(\mathbf{x},\hat{\mathbf{n}})$.  The framework applies identically to solids, liquids, gases, and individual particles: each is a State28D object that blocks CMB pressure along its zero-point line.


%----------------------------------------------------------------------
\section{Conclusion}
%----------------------------------------------------------------------

We have formalised the 28-dimensional aspect space as a minimal complete basis for displacement--spation interactions.  The occlusion function $E = T_3/(4\pi r^2)$ produces the $1/r^2$ force law from pure geometry.  The screening hierarchy ($10^{-9}$ to $10^{-123}$) resolves the force-hierarchy problem without new fields.  The movement principle ($d\Psi/dt \ge 0$) provides deterministic evolution, and three falsifiable experiments are proposed.

This chapter establishes the geometric foundation for all subsequent SDT dynamics.


% Bibliography references are consolidated in the main document.
